\section{数字孪生平台下的岩石识别框架}
\label{sec:ch5_dt_perception_framework}

为解决第\ref{sec:ch5_requirements}节提出的多域耦合难题，本文将岩石识别嵌入数字孪生闭环，构建“环境建模--感知特征提取--风险映射--规划反馈”的统一流程。
该流程在工程实现上对应环境建模模块、感知计算模块与运行调度模块的串联调用。

\begin{figure}[H]
	\centering
	\resizebox{0.94\textwidth}{!}{%
	\begin{tikzpicture}[
		node distance=6mm,
		box/.style={draw, rounded corners, align=left, minimum height=7mm, text width=0.82\textwidth, inner sep=2.5mm},
		arrow/.style={-{Stealth[length=2mm]}, thick, draw=black!75}
		]
		\node[box, fill=blue!10] (in) {\textbf{数据采集层：}DEM/DOM 先验 + 车载图像/点云 + 运行配置（分辨率、区域尺度、语义比例）};
		\node[box, fill=cyan!12, below=of in] (env) {\textbf{孪生环境构建层：}生成 $\{M_{elev},M_{sem},M_{phy},M_{trav}\}$，并形成统一环境数据包};
		\node[box, fill=green!12, below=of env] (perc) {\textbf{感知特征层：}基于高程梯度提取 $M_{slope}$ 与 $M_{rough}$，并保留语义掩膜};
		\node[box, fill=orange!12, below=of perc] (risk) {\textbf{风险映射层：}融合岩石占据、坡度与动力学约束，形成障碍风险场 $R_{obs}(x,y)$};
		\node[box, fill=purple!10, below=of risk] (plan) {\textbf{虚实交互层：}将风险场写入规划器；规划/动力学残差反向更新感知阈值与权重};
		\draw[arrow] (in) -- (env);
		\draw[arrow] (env) -- (perc);
		\draw[arrow] (perc) -- (risk);
		\draw[arrow] (risk) -- (plan);
		\draw[arrow, draw=blue!70!black] (plan.west) to[out=180,in=180,looseness=1.1] node[left,align=left]{在线修正\\(阈值/权重回写)} (perc.west);
	\end{tikzpicture}%
	}
	\caption{数字孪生平台下岩石识别框架与虚实闭环}
	\label{fig:ch5_dt_perception_framework}
\end{figure}

\subsection{数据采集与多层孪生建模}

环境构建阶段以统一栅格坐标组织四类基础图层：
\begin{equation}
\mathcal{M}_{DT}=\{M_{elev},M_{sem},M_{phy},M_{trav}\},
\label{eq:ch5_dt_map_set}
\end{equation}
其中 $M_{elev}$ 为高程图，$M_{sem}$ 为语义标签图，$M_{phy}$ 为物理参数图层（$k_c,k_\phi,n,c,\phi$），$M_{trav}$ 为可通行性图。

在代码实现中，语义标签与可通行性初始化规则为：松散月壤/压实月壤/岩石对应基准可通行性分别为 $0.7/0.9/0.4$，并叠加坡度惩罚项（式\eqref{eq:ch5_traversability}）。
该设计使感知层可直接继承环境物理先验，避免“视觉识别结果与动力学不可通行区域不一致”的系统偏差。

\subsection{感知特征提取与接口标准化}

项目当前轻量感知实现采用梯度特征快速提取策略：
\begin{equation}
M_{slope}(x,y)=\sqrt{\left(\frac{\partial z}{\partial x}\right)^2+\left(\frac{\partial z}{\partial y}\right)^2},
\quad
M_{rough}(x,y)=\left|\frac{\partial z}{\partial x}\right|+\left|\frac{\partial z}{\partial y}\right|.
\label{eq:ch5_slope_roughness}
\end{equation}

对应核心实现片段如下：
\begin{lstlisting}[language=Python,caption={感知特征提取核心代码片段},label={lst:ch5_perception_pipeline}]
data = np.load(env_artifact, allow_pickle=True)
elevation_map = data["elevation_map"]
semantic_map = data["semantic_map"]

grad_y, grad_x = np.gradient(elevation_map)
slope_map = np.sqrt(grad_x ** 2 + grad_y ** 2)
roughness_map = np.abs(grad_x) + np.abs(grad_y)
\end{lstlisting}

从系统集成角度，感知模块必须输出统一特征数据包并满足端到端调用契约：
\begin{equation}
\mathcal{I}_{per}=\{M_{slope},M_{rough},M_{sem},\mathcal{T}_{stamp},\mathcal{C}_{cfg}\},
\label{eq:ch5_perception_interface}
\end{equation}
其中 $\mathcal{T}_{stamp}$ 为时间戳，$\mathcal{C}_{cfg}$ 为配置快照。
该接口通过统一元数据记录机制保证实验可追溯性与跨模块一致性。

\subsection{虚实交互与在线更新机制}

感知层不是单向流水线，而是与规划和动力学形成反馈闭环。
设第 $k$ 时刻感知输出为 $\mathbf{z}_k$，规划反馈为 $\mathbf{u}_k$，动力学残差为 $\mathbf{e}_k$，则在线更新写为
\begin{equation}
\theta_{k+1}=\theta_k-\eta\nabla_{\theta}
\left(
\alpha\|\mathbf{z}_k-\hat{\mathbf{z}}_k\|_2^2+
\beta\|\mathbf{u}_k-\hat{\mathbf{u}}_k\|_2^2+
\gamma\|\mathbf{e}_k\|_2^2
\right),
\label{eq:ch5_online_update}
\end{equation}
其中 $\theta$ 为感知阈值/融合权重参数集，$\eta$ 为更新步长。

工程上，当前项目已实现“环境$\rightarrow$感知$\rightarrow$规划$\rightarrow$动力学”的端到端链路与统一结果归档。
在此基础上，后续第\ref{sec:ch5_orbslam}节将引入增强型 ORB-SLAM2 与 SiaT-Hough 语义模块，进一步提升岩石识别的几何一致性与鲁棒性。
